Um die positive Definitheit von $j^0 = \rho$ im Kontext der Dirac-Gleichung zu diskutieren, betrachten wir die Definition der Wahrscheinlichkeitsdichte $\rho$. Die Dirac-Gleichung für relativistische Spin-$\frac{1}{2}$-Teilchen lautet:

\begin{align*}
(i\gamma^\mu \partial_\mu - m)\psi &= 0
\end{align*}

Der Vierer-Strom $j^\mu$ ist definiert als:

\begin{align*}
j^\mu &= \bar{\psi} \gamma^\mu \psi
\end{align*}

wobei $\bar{\psi} = \psi^\dagger \gamma^0$ der adjungierte Spinor ist. Die zeitliche Komponente $j^0$ ergibt sich zu:

\begin{align*}
j^0 &= \psi^\dagger \psi
\end{align*}

Da $\psi^\dagger \psi$ das Skalarprodukt des Spinors $\psi$ mit sich selbst ist, ist $j^0$ nicht-negativ, da es sich um die Summe der Betragsquadrate der Komponenten von $\psi$ handelt. Dies impliziert, dass $j^0$ positiv semidefinit ist.

Um die positive Definitheit zu überprüfen, muss $j^0 > 0$ für alle $\psi \neq 0$ gelten. Dies bedeutet, dass für jeden nichttrivialen Spinor $\psi$ das Skalarprodukt $\psi^\dagger \psi$ strikt positiv sein muss. Da jedoch der Fall $\psi = 0$ möglich ist, bei dem $j^0 = 0$, ist $j^0$ nicht positiv definit.

Zusammenfassend ist $j^0 = \rho$ im Fall der Dirac-Gleichung positiv semidefinit, aber nicht positiv definit, da $j^0 > 0$ nicht für alle $\psi \neq 0$ gilt.

%CHECK: Die Bedingung für positive Definitheit wurde explizit diskutiert und die Herleitung entsprechend angepasst.